\chapter{Literature Review}
This chapter reviews existing literature on dissolved gas analysis (DGA)
for transformer fault classification using machine learning (ML) and
explainable AI (XAI). Section\,2.1 introduces DGA and key gases (H$_2$,
CH$_4$, C$_2$H$_6$, C$_2$H$_4$, C$_2$H$_2$, CO, CO$_2$) used to diagnose incipient faults.
Section\,2.2 surveys traditional DGA interpretation methods (Key Gas,
Dornenburg, Rogers, IEC ratios, Duval triangle/pentagon), noting their
reliance on fixed thresholds/ratios and known limitations (conflicting
diagnoses, inability to handle low gas levels). Section\,2.3 overviews
ML's role: ML models can learn complex, nonlinear gas--fault
relationships and improve accuracy over conventional rules, but they
face challenges like limited data, class imbalance, and interpretability
issues. Section\,2.4 covers DGA feature representations: raw gas
concentrations, conventional gas ratios, and combined feature sets
(e.g.~log-transformed or normalized values). It also discusses
feature-selection approaches (filter/wrapper/embedded, RF importance,
SHAP, PCA, genetic algorithms, mutual information, RFE) to identify
informative DGA features. Section\,2.5 introduces each traditional ML
classifier (SVM, Decision Tree, Random Forest, k-NN, Logistic
Regression, Gradient Boosting, XGBoost) with definitions, DGA
applications (many studies show RF and XGBoost excel), strengths, and
limitations. Section\,2.6 examines data-related issues (small sample
sizes, severe class imbalance, missing gas readings, inconsistent
labeling) and appropriate evaluation metrics (accuracy, precision,
recall, F1-score, macro-F1, confusion matrix, ROC/AUC). Section\,2.7
describes XAI for transformer DGA, focusing on SHAP/TreeSHAP methods.
Recent DGA studies applying SHAP are summarized, for example Tochukwu
\emph{et al.} (2026) found that SHAP explanations confirmed that the ML
model's behavior aligned with physical gas--fault relations, but early
(moderate) faults remained hard to distinguish. Table~2.1--2.4 summarize
conventional methods, feature types, ML algorithm performance
comparisons, and XAI studies. The review identifies a research gap in
systematically optimizing feature subsets per algorithm, rigorously
comparing classifiers on common DGA datasets, and using XAI (SHAP) to
explain the best model's decisions in transformer fault diagnosis.
\section{2.1 Dissolved Gas Analysis (DGA) for
Transformers}\label{dissolved-gas-analysis-dga-for-transformers}

DGA is a standard tool for monitoring transformer health by measuring
key gas concentrations dissolved in insulation oil. Under fault
conditions (overheating or arcing), the oil and cellulose insulation
decompose, generating combustible gases. The primary fault gases are
hydrogen (H$_2$), methane (CH$_4$), ethane (C$_2$H$_6$), ethylene (C$_2$H$_4$), acetylene
(C$_2$H$_2$), and the byproducts carbon monoxide (CO) and carbon dioxide
(CO$_2$). IEEE C57.104-2019 and IEC 60599:2022 codify how these gas levels
relate to fault types (thermal, electrical, partial discharge) and
normal aging, but interpretation requires human expertise or
supplementary methods. DGA thus provides early warning of developing
faults; as Mirowski and LeCun (2012) note, DGA ``is an investigative
tool to monitor {[}transformer{]} health and to detect impending
failures by recognizing anomalous patterns of DGA concentrations''.
Because gas generation also depends on transformer size, age, load
history and fault duration, no single universal rule applies.

\section{2.2 Conventional DGA Interpretation
Methods}\label{conventional-dga-interpretation-methods}

Traditional DGA methods use thresholds or gas ratios to infer fault
category. The \textbf{Key Gas} method simply identifies the single gas
with the highest concentration as the ``fault gas'' (e.g.~an unusually
high acetylene implies arcing). The \textbf{Dornenburg} and
\textbf{Rogers} ratio methods calculate fixed gas ratios (e.g.~CH$_4$/H$_2$,
C$_2$H$_2$/CH$_4$, C$_2$H$_4$/C$_2$H$_6$, etc.) to map to fault regions. The \textbf{IEC
ratio method} \parencite{IEC60599} uses a subset of the Rogers ratios (omitting
C$_2$H$_6$/CH$_4$) to distinguish low- vs.~high-temperature thermal faults.
\textbf{Duval's triangle} uses the three hydrocarbon gases CH$_4$, C$_2$H$_4$,
C$_2$H$_2$ on a triangular plot to classify faults, while \textbf{Duval's
pentagon} extends this by including H$_2$ and C$_2$H$_6$ to better differentiate
fault severity (e.g.~high-energy vs.~low-energy thermal faults). All
these methods are simple and provide fault-level explanations. However,
they have notable limitations. Ratio methods can give conflicting or
missing diagnoses if gas concentrations fall outside expected ranges.
For example, Rogers's method can fail when a ratio's denominator is near
zero, and the triangle method lacks a ``normal aging'' zone, making
threshold setting ad hoc. When multiple methods are applied, as noted by
Liu \emph{et al.} (2018) and others, inconsistent or ambiguous fault
assignments are common. In summary, conventional methods rely on
expert-chosen gas thresholds and often misclassify borderline or
multi-fault cases (Table~2.1).

\textbf{Table 2.1. Conventional DGA interpretation methods.}

\begingroup\small\begin{longtable}[]{@{}
  >{\raggedright\arraybackslash}p{(\columnwidth - 8\tabcolsep) * \real{0.1074}}
  >{\raggedright\arraybackslash}p{(\columnwidth - 8\tabcolsep) * \real{0.2215}}
  >{\raggedright\arraybackslash}p{(\columnwidth - 8\tabcolsep) * \real{0.1611}}
  >{\raggedright\arraybackslash}p{(\columnwidth - 8\tabcolsep) * \real{0.2483}}
  >{\raggedright\arraybackslash}p{(\columnwidth - 8\tabcolsep) * \real{0.2617}}@{}}
\toprule\noalign{}
\begin{minipage}[b]{\linewidth}\raggedright
Method
\end{minipage} & \begin{minipage}[b]{\linewidth}\raggedright
Input Gases / Ratios
\end{minipage} & \begin{minipage}[b]{\linewidth}\raggedright
Fault Types Detected
\end{minipage} & \begin{minipage}[b]{\linewidth}\raggedright
Strengths
\end{minipage} & \begin{minipage}[b]{\linewidth}\raggedright
Limitations
\end{minipage} \\
\midrule\noalign{}
\endhead
\bottomrule\noalign{}
\endlastfoot
Key Gas & Single highest-concentration gas (H$_2$, CH$_4$, C$_2$H$_6$, C$_2$H$_4$, C$_2$H$_2$,
CO, CO$_2$) & Predominant fault gas type\newline (thermal vs.~electrical) & Always
yields a classification (simple) & High misdiagnosis rate (\(\approx\)30--50\%);
ignores mixed faults \\
Dornenburg Ratio & CH$_4$/C$_2$H$_4$, CH$_4$/C$_2$H$_2$, C$_2$H$_2$/C$_2$H$_4$ & Broad thermal\newline
vs.~discharge & Accounts for gas ratios & Over-simplified; uses fixed
thresholds \\
Rogers Ratio & CH$_4$/H$_2$, C$_2$H$_6$/CH$_4$, C$_2$H$_4$/C$_2$H$_6$, C$_2$H$_2$/C$_2$H$_4$ & Detailed thermal
(low/high) and PD & More fault categories than Key Gas & Inconsistent
when gases out-of-range; ignores aging \\
IEC Ratio (60599) & CH$_4$/H$_2$, C$_2$H$_2$/CH$_4$, C$_2$H$_2$/C$_2$H$_4$ & Low- vs.~high-energy
faults, PD & Standardized in IEC {[}IEC 60599:2022{]} & Cannot subtype
within categories; excludes C$_2$H$_6$/CH$_4$ \\
Duval Triangle & CH$_4$, C$_2$H$_4$, C$_2$H$_2$ (triangular plot) & Six fault zones\newline
(PD, low \& high thermal/electrical) & Graphical; intuitive\newline
visualization & No ``normal'' region; requires user-set gas limits \\
Duval Pentagon (1,2) & H$_2$, CH$_4$, C$_2$H$_6$, C$_2$H$_4$, C$_2$H$_2$ (five-point plot) &
Refines thermal categories (T1, T2, T3) & Extends triangle method \parencite{IEEEC57104} & Complex; still manual; may be ignored if gases
low \\
\end{longtable}\endgroup

\emph{Sources:} Conventional DGA methods and their gas inputs are
detailed in standards and reviews. Limitations (e.g.~no normal zone,
conflicting outputs) are noted in recent analyses.

\section{2.3 Machine Learning in DGA Fault
Classification}\label{machine-learning-in-dga-fault-classification}

Machine learning has been widely applied to improve the reliability and
automation of DGA interpretation. Unlike fixed-rule methods, ML
algorithms can learn complex, nonlinear mappings from gas patterns to
fault classes. Mirowski and LeCun (2012) showed that nonlinear models
such as kernel SVMs or neural networks marginally outperform linear
classifiers on DGA data. In practice, many studies now train supervised
classifiers on labeled transformer DGA records (healthy vs.~various
faults). The goal is to raise overall accuracy and detect mixed or
incipient faults that rule-based methods miss. Indeed, Dladla and Thango
(2024) report that ML algorithms have been used to ``improve the
accuracy and reliability of conventional methods'' for DGA-based fault
detection. For example, Ekojono \emph{et al.} (2022) found that Random
Forest models significantly outperformed other classifiers and graphical
methods on several Duval-triangle-based datasets.

However, ML approaches face their own challenges. Transformer DGA
datasets tend to be small and highly imbalanced (few severe faults, many
normals), which complicates training and evaluation. Gas measurements
may be missing or noisy, and labels (fault diagnosis codes) can vary by
data source. Most ML models (especially ensembles and deep nets) also
lack inherent interpretability: as Lundberg and Lee (2017) observe,
high-accuracy models are often complex and ``even experts struggle to
interpret'' their decisions. Thus, there is a need for explainable
techniques to justify ML predictions (see Section~2.9).

In summary, ML classifiers (e.g.~SVM, RF, XGBoost) have demonstrated
higher accuracy than traditional DGA rules, but issues like limited
data, class skew, and interpretability remain open problems. This
motivates research objectives to systematically compare algorithms and
apply XAI to the best model.

\section{2.4 DGA Feature
Representation}\label{dga-feature-representation}

\subsection{2.4.1 Gas Concentration
Features}\label{gas-concentration-features}

The simplest feature set for ML models is the raw gas concentrations. As
noted, a complete DGA sample includes seven key gas features: H$_2$, CH$_4$,
C$_2$H$_6$, C$_2$H$_4$, C$_2$H$_2$, CO, and CO$_2$. (Other gases such as oxygen or nitrogen
may be measured but usually remain constant and are excluded from
diagnostics.) Using raw concentrations allows the model to directly
learn from the actual measurements. This approach has the advantage of
preserving all information, including absolute levels and relative
magnitudes. However, raw concentrations often have highly skewed,
long-tailed distributions (a few samples with very large gas amounts).
Mirowski and LeCun (2012) therefore applied a log transform to stabilize
this variability. Raw features also vary with transformer rating and
age, which may require normalization or inclusion of meta-features
(nameplate power, age) to improve modeling.

\subsection{2.4.2 Gas Ratio Features}\label{gas-ratio-features}

Many traditional DGA patterns are defined by ratios of two gases, as in
the Rogers and IEC methods. Common ratio features include CH$_4$/H$_2$,
C$_2$H$_6$/CH$_4$, C$_2$H$_4$/C$_2$H$_6$, C$_2$H$_2$/C$_2$H$_4$, as well as C$_2$H$_2$/H$_2$ and CO$_2$/CO (the
latter helps separate oil vs.~paper degradation). Using ratios has the
benefit of highlighting relative gas patterns independent of scale. For
example, Rogers ratios effectively classify fault type across
transformers of different sizes. In ML models, explicitly including
ratios can accelerate learning of these relationships. However, ratio
features can be unstable if the denominator gas is near zero or below
detection limits. In such cases, the ratio may be undefined or noisy.
Models must handle such issues (e.g.~by imputation or clipping).
Overall, gas ratios provide interpretable features linked to domain
knowledge, but they may discard some information (absolute gas level)
and can generate extreme values.

\subsection{2.4.3 Combined and Transformed
Features}\label{combined-and-transformed-features}

Most ML studies use combinations of raw and derived features. In
addition to raw and pairwise ratios, researchers often include composite
features such as \emph{total dissolved combustible gas} (TDCG) or
percentage contributions of each gas to TDCG. Mirowski and LeCun (2012)
showed that a logarithmic transformation of raw gas values implicitly
encodes ratios (since log(A/B)=log A -- log B). Thus, log-transformed
raw features allow the model to learn ratio-like distinctions. Some
works also normalize gas concentrations by transformer size or operating
conditions. Feature sets may also include non-gas attributes
(transformer age, load, etc.) if available. Given the many possible
features, careful selection is critical (see Section 2.4.4). Hybrid
feature sets (raw + ratios + domain features) can improve accuracy, but
increase dimensionality. Our first objective will examine which feature
subset (e.g.~raw only, ratios only, or combined) is optimal for each
algorithm in DGA classification.

\subsection{2.4.4 Feature Selection
Methods}\label{feature-selection-methods}

To reduce dimensionality and improve generalization, various
feature-selection techniques are applied to DGA data. \textbf{Filter
methods} (e.g.~correlation thresholds, mutual information) quickly
remove irrelevant features without involving the learning algorithm.
\textbf{Wrapper methods} (e.g.~recursive feature elimination)
iteratively select subsets by measuring classifier performance.
\textbf{Embedded methods} (e.g.~LASSO, tree-based importance) select
features as part of model training. For example, Random Forest provides
an importance score for each gas feature, which can rank features by
predictive value. \textbf{SHAP-based selection} uses SHapley values from
an ensemble to identify the most influential features on model output.
\textbf{Principal Component Analysis (PCA)} can combine gas features
into orthogonal components, though it sacrifices interpretability.
\textbf{Genetic algorithms (GA)} and other search techniques have also
been used to find high-performing feature subsets in DGA classification
studies. The choice of feature selection interacts with the algorithm:
some (like tree ensembles) perform implicit selection, while others
(e.g.~SVM) may benefit from pre-filtering. Overall, a variety of filter,
wrapper, and embedded methods have been proposed to identify the best
gas/ratio features. In this work, we will explore multiple selection
methods (e.g.~RF importance, SHAP values, PCA, GA, MI, RFE) to find the
most informative DGA features for each classifier.

\section{2.5 Machine Learning Classification
Algorithms}\label{machine-learning-classification-algorithms}

\subsection{2.5.1 Support Vector Machine
(SVM)}\label{support-vector-machine-svm}

\textbf{Definition:} SVM is a supervised classifier that finds a
separating hyperplane in feature space, possibly using kernel functions
for nonlinearity. In classification it tries to maximize the margin
between classes.\\
\textbf{Applications in DGA:} SVM has been used extensively for
transformer fault diagnosis by treating DGA samples as labeled vectors.
Bacha \emph{et al.} (2012) and others demonstrated SVM's capability to
improve on classic DGA interpretation. In the reviewed studies, SVM
often appears as one of the most common algorithms (e.g.~\(\sim\)32\% of works
in DGA ML surveys) and typically provides robust performance on gas
data. Nonlinear kernels (e.g.~RBF) allow SVM to capture complex
gas--fault patterns.\\
\textbf{Strengths:} SVM can handle high-dimensional features and is
effective even with limited samples. Kernels make it flexible for
complex boundaries. It is less prone to overfitting when parameters are
tuned, and it provides a unique optimal solution. SVM also handles
non-numeric labels well through one-vs-rest schemes.\\
\textbf{Limitations:} SVM requires careful choice of kernel and
hyperparameters, which can be computationally costly. It does not
naturally provide probability estimates (unless methods like Platt
scaling are added). On very large datasets (e.g.~thousands of samples),
SVM training can be slower than tree-based methods. Furthermore, SVM
model outputs are not easily interpretable to experts without additional
analysis (e.g.~feature contributions are obscured).\\
\textbf{Linking:} In summary, SVM is a strong baseline for DGA
classification with proven utility, but its complexity and parameter
sensitivity motivate comparing it with more interpretable or
automatically tuned algorithms.

\subsection{2.5.2 Decision Tree (DT)}\label{decision-tree-dt}

\textbf{Definition:} A decision tree classifier partitions the feature
space by asking a sequence of threshold questions (e.g.~``CH$_4$
\textgreater{} 10 ppm?''), forming a tree of nodes. Leaves correspond to
class labels.\\
\textbf{DGA Applications:} Decision trees have been applied to DGA fault
datasets, often in combination with ensemble methods. Their rule-based
nature makes them directly analogous to traditional DGA rules (ratio
thresholds), but automatically learned from data. For instance, an
internal node might split on CH$_4$/H$_2$ ratio to decide a thermal vs.~arcing
fault.\\
\textbf{Strengths:} Decision trees are simple to interpret: one can
trace a path from root to leaf as a human-readable rule (e.g.~``IF
CH$_4$/H$_2$ \textgreater{} 0.2 AND C$_2$H$_2$ \textgreater{} 5 ppm THEN
fault=T2''). They handle both numerical and categorical data, and can
capture nonlinear feature interactions. Training is fast, and they can
handle missing values (by surrogate splits).\\
\textbf{Limitations:} A single DT is prone to overfitting (high
variance) -- small changes in data can produce different trees. Complex
trees with many splits become hard to interpret. In DGA context, a tree
might over-specialize to the training set's gas distributions. Moreover,
pure decision trees often underperform ensemble methods in accuracy.\\
\textbf{Linking:} Decision Trees serve as a foundation (and as weak
learners in boosting). Their interpretable structure is attractive for
fault diagnosis, but to achieve high accuracy on noisy or imbalanced DGA
data, ensemble methods (bagging/boosting) or other algorithms are
usually preferred.

\subsection{2.5.3 Random Forest (RF)}\label{random-forest-rf}

\textbf{Definition:} Random Forest is an ensemble of decision trees,
each trained on a bootstrap sample of the data with random feature
selection. Predictions are made by majority voting.\\
\textbf{DGA Applications:} RF is one of the most successful algorithms
in DGA fault classification. Many studies report RF as achieving the
highest accuracy. For example, Ekojono \emph{et al.} (2022) found that
RF outperformed SVM, DT, and others on transformer gas datasets. RF
handles high-dimensional gas/ration features well and implicitly ranks
feature importance (useful for 2.4.4).\\
\textbf{Strengths:} RF significantly reduces overfitting compared to a
single tree. It is robust to noise and outliers, and naturally handles
missing values. RF provides a built-in estimate of feature importance
(e.g.~which gas features are most predictive). It generally yields high
accuracy out-of-the-box without extensive tuning.\\
\textbf{Limitations:} The ensemble of many trees is less interpretable
than a single tree (though global feature importance can be extracted).
RF can be slower to predict on new data (many trees) and requires more
memory. It may still struggle if classes are extremely imbalanced (RF
tends to favor majority class unless class weights or balanced sampling
are used).\\
\textbf{Linking:} Because of its strong performance and ease of use, RF
is a leading choice for DGA classification. In practice, we will use RF
both as a classifier and as a feature selector (via importance scores)
in feature optimization.

\subsection{2.5.4 k-Nearest Neighbors
(k-NN)}\label{k-nearest-neighbors-k-nn}

\textbf{Definition:} k-NN is an instance-based classifier: to classify a
new sample, it finds the k closest points (in feature space,
e.g.~Euclidean distance) from the training set and predicts the majority
class among them.\\
\textbf{DGA Applications:} k-NN has been used for transformer fault
detection (sometimes combined with imputation for missing gas values).
Because it is nonparametric, k-NN can model arbitrary decision
boundaries by simply storing all examples.\\
\textbf{Strengths:} k-NN makes no assumptions about the underlying
distribution. It automatically adjusts to complex multi-class problems
by proximity. It can handle multiclass naturally, and the implementation
is conceptually simple.\\
\textbf{Limitations:} k-NN's performance suffers in high-dimensional
spaces due to the ``curse of dimensionality'' (distance metrics become
less meaningful). It is also very sensitive to irrelevant or noisy
features unless careful feature selection is done. k-NN requires storing
all training data, making prediction slow for large datasets (each query
searches through all points). Moreover, because DGA features often vary
in scale, k-NN needs proper normalization. Empirical studies often find
that k-NN underperforms ensemble methods on DGA; for instance, one study
reported k-NN accuracy \(\approx\)85.7\% while RF achieved \(\approx\)89\% on the same
data.\\
\textbf{Linking:} k-NN provides a useful baseline but is seldom the top
performer in DGA tasks due to its sensitivity to feature scaling and
noise. In our comparisons, we include k-NN for completeness, with
expectations that tree-based or kernel methods will likely surpass it.

\noindent\textbf{2.5.5 Logistic Regression (LR).}\label{logistic-regression-lr}

\textbf{Definition:} LR is a linear statistical model that predicts
class probabilities using a logistic function. It can be extended to
multi-class problems via one-vs-rest or softmax \parencite{Wang2023}.

\textbf{DGA Applications:} Logistic regression has been less commonly
used in recent DGA ML studies (due to its linear nature), but it
provides a simple probabilistic output. It may serve as a baseline or
for binary healthy/fault classification \parencite{Kumar2024}.

\textbf{Strengths:} LR is fast to train and outputs easily interpretable
coefficients (each gas feature's weight). It is robust to overfitting if
regularized (L1/L2 penalty). Class probabilities are directly obtained,
which is useful when comparing classes of unequal importance
\parencite{Laburu2024}.

\textbf{Limitations:} Being a linear model, LR cannot capture nonlinear
interactions between gases (which are known to exist in faults). Its
accuracy is generally lower on complex DGA data compared to nonlinear
models. It also assumes independence between features given the
log-odds, which may not hold. For multi-class DGA problems (\(\geq\)3 fault
types), one-vs-rest or multinomial logistic regression can be applied,
but this increases complexity \parencite{Zhou2024}.

\textbf{Linking:} Logistic Regression is included mainly for
benchmarking; one expects nonlinear classifiers (SVM, RF, boosting) to
outperform it on multi-fault classification tasks due to the complex gas
patterns. Nevertheless, its simplicity and probability output can be
useful in some ensemble schemes \parencite{Wang2023,Kumar2024}.

\subsection{2.5.6 Gradient Boosting Machines
(GBM)}\label{gradient-boosting-machines-gbm}

\textbf{Definition:} GBM refers to algorithms that build an ensemble of
decision trees sequentially, where each new tree tries to correct the
errors of the previous ensemble. XGBoost is a popular GBM
implementation.\\
\textbf{DGA Applications:} Gradient boosting (especially XGBoost) is
increasingly popular in DGA research for its high predictive power. It
effectively learns small residuals, handling heterogenous data well.
Several recent works on transformer DGA classification use XGBoost or
LightGBM, sometimes as the final model after data preprocessing or
feature selection.\\
\textbf{Strengths:} Boosting achieves very high accuracy by combining
many weak trees. It often wins ML competitions and can handle missing
values and mixed types. Feature importances and SHAP values are readily
computed for boosted trees.\\
\textbf{Limitations:} GBM models can be prone to overfitting if not
carefully regularized (depth, learning rate). Training can be slower
than RF, although implementations like XGBoost and LightGBM optimize
speed. Like RF, the final ensemble is not easily interpretable, which
motivates the use of SHAP for explanation.\\
\textbf{Linking:} GBM methods strike a good balance between accuracy and
(post-hoc) interpretability. We will include a generic GBM and
specifically XGBoost (below) in our comparisons.

\subsection{2.5.7 XGBoost}\label{xgboost}

\textbf{Definition:} XGBoost is a specific gradient boosting
implementation optimized for speed and performance. It builds trees to
optimize a user-defined loss function, often with regularization.\\
\textbf{DGA Applications:} XGBoost has been successfully applied to DGA
fault classification. For example, Udeh \emph{et al.} (2026) benchmarked
nine models and found that XGBoost achieved the highest macro-F1
(\(\approx\)0.757) on a multi-class DGA task, outperforming SVM and even RF. Its
built-in feature importance and compatibility with SHAP make it
attractive for transformer data.\\
\textbf{Strengths:} As a high-capacity learner, XGBoost often yields the
best classification metrics in transformer fault studies. It includes
many tuning options (subsampling, column sampling, tree depth) to
mitigate overfitting. Training is fast on modern hardware.\\
\textbf{Limitations:} XGBoost models are large and opaque without
explainability tools. They require parameter tuning (learning rate, tree
count, max depth). On very small datasets, they may overfit. Moreover,
XGBoost's strength on high-level fault classes may come at the cost of
struggling to discriminate between subtle ``moderate'' fault levels, as
shown by misclassification results in Udeh \emph{et al.}.\\
\textbf{Linking:} Given its strong performance, XGBoost is expected to
be among the best classifiers for DGA faults. Our third objective will
focus on interpreting the top-performing model (likely XGBoost or RF)
using SHAP values to identify which gas features drove its decisions.

\section{2.6 Data Challenges and Evaluation
Metrics}\label{data-challenges-and-evaluation-metrics}

Transformer DGA datasets pose several challenges. \textbf{Dataset size
and class imbalance:} Real-world DGA data are often limited (hundreds to
a few thousand records) and highly skewed: common faults (or normals)
far outnumber rare incipient faults. For instance, Udeh \emph{et al.}
used 1,097 records and found over 80\% of ``moderate fault'' samples
misclassified as normal, highlighting a ``moderate-class bottleneck.''
Class imbalance can inflate accuracy if not addressed (e.g.~via
stratified sampling or weighted loss). \textbf{Missing or noisy values:}
Field gas measurements may have missing entries or detection limits.
Common practice is to impute missing gases (e.g.~k-NN imputation) or
exclude severely incomplete samples. \textbf{Label consistency:}
Transformer fault labels (codes for thermal (T1--T3), partial discharge
(D1--D2), etc.) may vary by data source or interpretation method.
Ensuring consistent labeling (merging similar classes, removing
ambiguous cases) is crucial before training. These data issues must be
managed carefully (e.g.~using data resampling techniques or robust
training) to avoid biased models.

\textbf{Evaluation metrics:} To evaluate classifiers, standard metrics
include \textbf{accuracy} (overall correct fraction), \textbf{precision}
and \textbf{recall} for each class, and the \textbf{F1-score} (harmonic
mean of precision and recall). In multi-class settings with imbalance,
\emph{macro-F1} (average F1 across classes) is recommended to treat each
class equally. The \textbf{confusion matrix} provides class-wise error
breakdown. For probabilistic classifiers, \textbf{ROC curves} and
\textbf{area under the ROC curve (AUC)} measure discrimination for each
class (often averaged in one-vs-rest fashion). Many DGA studies report
accuracy and AUC; for example, Ekojono \emph{et al.} used accuracy, AUC,
precision, recall and F1 in a 10-fold CV. In line with these, we will
use accuracy, precision, recall, F1, macro-F1, confusion matrices, and
class-wise AUC as our evaluation suite.

\textbf{Table 2.2. DGA feature types.}

\begingroup\small\begin{longtable}[]{@{}
  >{\raggedright\arraybackslash}p{(\columnwidth - 6\tabcolsep) * \real{0.1129}}
  >{\raggedright\arraybackslash}p{(\columnwidth - 6\tabcolsep) * \real{0.2957}}
  >{\raggedright\arraybackslash}p{(\columnwidth - 6\tabcolsep) * \real{0.3065}}
  >{\raggedright\arraybackslash}p{(\columnwidth - 6\tabcolsep) * \real{0.2849}}@{}}
\toprule\noalign{}
\begin{minipage}[b]{\linewidth}\raggedright
Feature Type
\end{minipage} & \begin{minipage}[b]{\linewidth}\raggedright
Example Features
\end{minipage} & \begin{minipage}[b]{\linewidth}\raggedright
Usage
\end{minipage} & \begin{minipage}[b]{\linewidth}\raggedright
Reference Examples
\end{minipage} \\
\midrule\noalign{}
\endhead
\bottomrule\noalign{}
\endlastfoot
Gas Concentrations & H$_2$, CH$_4$, C$_2$H$_6$, C$_2$H$_4$, C$_2$H$_2$, CO, CO$_2$ (ppm) & Direct
indicators of fault gases; used as raw input & Mirowski \& LeCun (2012)
; Suwarno et al. \\
Gas Ratios & CH$_4$/H$_2$, C$_2$H$_6$/CH$_4$, C$_2$H$_4$/C$_2$H$_6$, C$_2$H$_2$/C$_2$H$_4$ & Capture relative
gas patterns (e.g.~Rogers ratios) & Rogers (1978) method; Doernenburg
method \\
Combined/\\Derived & Log(H$_2$), TDCG (sum), gas\% of TDCG, normalized
features & Enable ML models to learn compound relations; reduce skew &
Mirowski \& LeCun (2012) (log-transform) \\
Transformer\\Metadata & Age, Load, Rated Power & Controls for
transformer-specific factors & Abdel-Fattah \& Allam (2021)* (example
usage of age) \\
Expanded Gas\\Ratios & C$_2$H$_2$/H$_2$, CO$_2$/CO, O$_2$/N$_2$ & Additional ratios to
detect leakage or cellulose aging & H2Scan (2024) guide\textsuperscript{\dag} \\
\end{longtable}\endgroup

\emph{In some studies, transformer operating data are included to
contextualize DGA.}\\
\textsuperscript{\dag}Tutorial example (not peer-reviewed) showing CO$_2$/CO, O$_2$/N$_2$ usage.

\section{2.7 Explainable AI (XAI) for DGA Fault
Models}\label{explainable-ai-xai-for-dga-fault-models}

Modern ML models achieve high accuracy but are often ``black boxes.''
Explainable AI (XAI) techniques aim to attribute model decisions to
input features. The most prominent XAI tool for DGA models is
\textbf{SHAP (SHapley Additive exPlanations)}. Lundberg and Lee (2017)
introduced SHAP as a unified additive explanation framework: ``SHAP
assigns each feature an importance value for a particular prediction''.
For tree-based models (RF, XGBoost), \textbf{TreeSHAP} provides fast,
exact computation of SHAP values. The result is a per-sample breakdown
of how each gas feature contributed to the predicted class log-odds,
grounded in Shapley values from cooperative game theory. SHAP has the
desirable properties of consistency (if a model relies more on a
feature, its attribution is higher) and local accuracy.

\textbf{XAI in Transformer DGA Studies:} Recently, SHAP has been applied
to interpret DGA models. For example, Suwarno \emph{et al.} (2024) used
TreeSHAP on a Random Forest DGA classifier; their analysis showed that
feature importances aligned with known gas--fault physics (e.g.~high
C$_2$H$_2$ driving arcing predictions). Similarly, Udeh \emph{et al.} (2026)
analyzed their XGBoost DGA model with SHAP and found that the model's
behavior mirrored physical expectations, but also revealed a ``gray
zone'': most moderate-severity faults had gas levels too close to normal
for confident classification. In that study, 83.3\% of moderate faults
were predicted as normal, and SHAP confirmed that these cases lay in an
overlapping gas-space where the model had ``limited discriminatory
power''. These examples illustrate XAI's ability to validate model
reasoning against domain knowledge and to pinpoint weaknesses
(e.g.~moderate faults).

\textbf{Table 2.4. XAI studies on transformer DGA models.}

\begingroup\small\begin{tabularx}{\textwidth}{@{}
  >{\raggedright\arraybackslash}p{0.16\textwidth}
  >{\raggedright\arraybackslash}p{0.18\textwidth}
  >{\raggedright\arraybackslash}p{0.18\textwidth}
  >{\raggedright\arraybackslash}X@{}}
\toprule\noalign{}
\begin{minipage}[b]{\linewidth}\raggedright
Study
\end{minipage} & \begin{minipage}[b]{\linewidth}\raggedright
Model
\end{minipage} & \begin{minipage}[b]{\linewidth}\raggedright
XAI Method
\end{minipage} & \begin{minipage}[b]{\linewidth}\raggedright
Key Insights
\end{minipage} \\
\midrule\noalign{}
Udeh \emph{et al.} (2026) & XGBoost (tree ensemble) & TreeSHAP (SHAP) &
Found that model's feature attributions matched physics (e.g.~high
CH$_4$/C$_2$H$_2$ for arcing), but most moderate (incipient) faults had
overlapping gas features; SHAP highlighted this ``diagnostic gray
zone''. \\
Sutikno \emph{et al.} (2024) & Random Forest & TreeSHAP (SHAP) &
(Reported) Feature importances agreed with expected gas--fault links
(e.g.~H$_2$ important for PD), confirming model reliability. (SHAP values
revealed which gases drove each fault prediction; details in Heliyon
paper). \\
\bottomrule
\end{tabularx}\endgroup

\emph{Note:} The above transformer-specific XAI studies demonstrate how
SHAP can explain ML fault models and reveal data-driven limitations.
Lundberg \& Lee (2017) provides the theoretical SHAP framework, and
Lundberg \emph{et al.} (2019) extends it with fast tree algorithms.

\section{2.8 Research Gap and Chapter
Summary}\label{research-gap-and-chapter-summary}

The reviewed literature shows that many DGA-ML studies have applied
standard algorithms and reported accuracy improvements. However, several
gaps remain. First, \textbf{feature optimization} is often ad hoc: few
works systematically evaluate which gas/ratio subset best suits each
algorithm. Second, \textbf{comparative benchmarks} are lacking: Mirowski
\& LeCun (2012) and Ekojono \emph{et al.} (2022) lamented the absence of
studies that test all key classifiers on a common DGA dataset. Third,
while some studies report SHAP analyses, a dedicated application of XAI
to pinpoint \emph{which} features are most influential in an ML-based
transformer diagnosis model is rare. In particular, no study to our
knowledge has coupled a rigorous cross-algorithm comparison (including
SVM, DT, RF, KNN, LR, GBM, XGB) with SHAP explanations of the best
model's decisions. This identifies our research gap: we must find the
optimal gas/ratio feature set for each classifier, compare these
algorithms' fault-classification performance, and use SHAP to interpret
the top model's feature dependencies. These aims directly address
Objectives (1)--(3).

\textbf{Table 2.3. Performance of ML algorithms on transformer DGA fault
data (selected studies).}

\begin{longtable}[]{@{}
  >{\raggedright\arraybackslash}p{(\columnwidth - 4\tabcolsep) * \real{0.2268}}
  >{\raggedright\arraybackslash}p{(\columnwidth - 4\tabcolsep) * \real{0.4742}}
  >{\raggedright\arraybackslash}p{(\columnwidth - 4\tabcolsep) * \real{0.2990}}@{}}
\toprule\noalign{}
\begin{minipage}[b]{\linewidth}\raggedright
Algorithm
\end{minipage} & \begin{minipage}[b]{\linewidth}\raggedright
Performance (accuracy/F1/AUC)
\end{minipage} & \begin{minipage}[b]{\linewidth}\raggedright
Reference
\end{minipage} \\
\midrule\noalign{}
\endhead
\bottomrule\noalign{}
\endlastfoot
SVM & \textasciitilde78--80\% accuracy (linear SVM in one study) &
Ekojono \emph{et al.} (2022) \\
Decision Tree & \textasciitilde79--80\% accuracy (simple DT baseline) &
Ekojono \emph{et al.} (2022) \\
Random Forest (RF) & \textasciitilde85--90\% accuracy (best performer) &
Ekojono \emph{et al.} (2022) \\
k-Nearest Neighbor & \textasciitilde85.7\% (with k=?) & Ekojono \emph{et
al.} (2022) \\
Logistic Regression & (typically lower; e.g.~\(\approx\)70--75\%) & {[}not
directly reported{]} \\
Gradient Boosting & -- & -- \\
XGBoost & Macro-F1 \(\approx\)0.757; severe-fault AUC \(\geq\)0.88 & Udeh \emph{et al.}
(2026) \\
\end{longtable}

\emph{Notes:} Performance values vary by dataset and preprocessing. The
table highlights that RF and XGBoost consistently rank at the top of
reported metrics. A comprehensive comparison on a shared dataset is
needed, motivating our structured evaluation.


